\begin{env}[6.15.4]
\EGAReaderTarget{ega4.c04.statement.6-15-4}{IV.2.s6.15.4-fr}
Generalizando la definición dada en (\textbf I, 2.2.9), diremos que,
para dos preesquemas \emph{reducidos} $X$ e $Y$, un morfismo
$f:X\to Y$ es \emph{birracional} si la restricción de $f$ al
conjunto de los puntos maximales de $X$ es una \emph{biyección} de
este conjunto sobre el conjunto de los puntos maximales de $Y$
y si, para todo punto maximal $y$ de $Y$, el morfismo
$X\times_Y\operatorname{Spec}(\mathcal O_y)\to\operatorname{Spec}(\mathcal O_y)$
inducido por $f$ es un isomorfismo (es decir, si la fibra
$f^{-1}(y)$ se reduce a un solo punto $x$ (maximal en $X$)
y si el homomorfismo $\boldsymbol{k}(y)\to \boldsymbol{k}(x)$ inducido por $f$ es
biyectivo). Esta noción se reduce a la de (\textbf I, 2.2.9) cuando
$X$ e $Y$ solo tienen un número \emph{finito} de componentes
irreducibles.
\end{env}

\begin{lemma}[6.15.4.1]
\label{IV.2.s6.15.4.1-fr}
Sean $f:X\to Y$ un morfismo y $g:Y'\to Y$ un morfismo plano;
sean $X'=X\times_Y Y'$ y $f'=f_{(Y')}:X'\to Y'$. Si
$f_{\mathrm{red}}$ es birracional, también lo es
$f'_{\mathrm{red}}$.
\end{lemma}

\begin{proof}
En efecto, si $y'$ es un punto maximal de $Y'$, se sabe
(2.3.4, (ii)) que $y=g(y')$ es un punto maximal de $Y$.
Por hipótesis, existe un único punto $x$ de $X$ en
$f^{-1}(y)$; además, $x$ es maximal en $X$ y
$\boldsymbol{k}(x)=\boldsymbol{k}(y)$. De (\textbf I, 3.4.9) resulta que $f'^{-1}(y')$ se
reduce a un único punto $x'$ y que $\boldsymbol{k}(x')=\boldsymbol{k}(y')$.
Se concluye observando que, por (2.3.7, (ii)), toda
componente irreducible de $X'$ domina una componente
irreducible de $Y'$.
\end{proof}
